\subsection{Analytic radial entrance for the actual corrected central field}
\label{sec:cq-actual-angular-analytic}

The preceding general entrance criteria can be applied to the actual
central field without an additional analyticity assumption. We retain
the profiles and all cutoff choices of the cited source. The argument
below proves the required radial analyticity locally at the axis; it
does not claim that a general smooth cutoff is analytic on its
transition region.

First consider the original viscosity-one similarity variables
\[
 \tau=1-t,\qquad \tau=q(1-\eta^2),\qquad
 z=q^{1/2-h}\eta,\qquad X=\frac{r^2}{2q}.
\]
For every preterminal point in the source's local domain, $q>0$.
The source's Proposition B.2 supplies a leading profile jointly
analytic in $X$ and $\eta$ on a neighborhood of its initial radial
rectangle. Every positive-order profile actually selected by the
source is analytic in $X$ near $X=0$, as we now prove from its
displayed recurrence.

Set $\xi=\sqrt X$. Source equations (5.2)--(5.7) give at each
positive order $n$ the linear regular-singular system
\begin{equation}
\begin{split}
 \partial_\xi W_n+\xi^{-1}C_{\rm diag}W_n
       &=A_0(\xi,\eta)W_n+
             A_1(\xi,\eta)\partial_\eta W_n+f_n(\xi,\eta),\\
 C_{\rm diag}&=\operatorname{diag}(0,0,2,0,3,1),\\
 W_n&=(\phi_n,U_n,K_n,\Pi_n,\partial_\xi\phi_n,
                                      \partial_\xi U_n)^{\mathsf T},\\
 K_n&=A_X(U_n)-U_n,\qquad W_n(0,\eta)=0 .
\end{split}
 \label{eq:cq-actual-axis-linear-system}
\end{equation}
The matrices and source here are the unchanged ones from that
recurrence. In particular the only possibly nonzero entries of
$A_1$ are exactly
\begin{align*}
 (A_1)_{51}&=2H_c/L,&
 (A_1)_{52}=(A_1)_{53}
     &=-2d(\phi_0+X\partial_X\phi_0)/L,\\
 (A_1)_{62}&=2(H_c-dX\partial_XU_0)/L,&
 (A_1)_{63}&=-2dX\partial_XU_0/L,&
 (A_1)_{64}&=2d/L,
\end{align*}
where $d=1-\eta^2$, $L=1-2h\eta^2$ and
$H_c=(1/2-h)\eta+dU_0$. Thus $A_1$ and every parameter derivative
map the first four components to the last two and annihilate
the last two.

Assume inductively that all lower-order profiles are analytic near
the axis. This is true at order zero by the just-cited proposition.
The source's radial extensions agree with their inner profiles on
one fixed initial radial interval; their moment corrections have
support to the right of that interval. Thus those extensions do not
change any lower-order germ used here. The equations forming
$A_0,A_1,f_n$ use finitely many sums, products, radial and parameter
derivatives, and the original radial averages, with divisions only
by $L$, nonzero leading factors on the chosen neighborhood, and
the powers of $X$ whose zero factors are explicitly canceled in
the source recurrence. These operations preserve the analytic
germs. For the average the exact formula is
$A_Xg=\int_0^1g(bX,\eta)\,db$. For the apparently divided sources,
$V_j=Xv_j$ by (5.2), and every summand of $\Omega_j$ in (5.6)
is divisible by $X$, including
\[
 V_i(\partial_XV_j-V_j/(2X))
       =Xv_i(v_j/2+X\partial_Xv_j).
\]
Analytic divisibility follows from the vanishing Taylor coefficient,
so no singular coefficient is discarded. Once the displayed
diagonal singular part is separated, $A_0,A_1,f_n$ consequently
extend holomorphically to a product
$|\xi|<R$, $\eta\in\Omega$ about the chosen real axis point.
Their norms are bounded on a smaller closed product. The radii
may depend on $n$; only positivity of each radius is used below.

For $c_i=(C_{\rm diag})_{ii}$ extend the exact source inverse
to complex $\xi$ by
\[
 (Gg)_i(\xi,\eta)=
       \xi\int_0^1b^{c_i}g_i(b\xi,\eta)\,db,\qquad
 K=G(A_0+A_1\partial_\eta).
\]
It is holomorphic on the product, has value zero at $\xi=0$,
and termwise integration of a convergent Taylor series proves
$(\partial_\xi+c_i/\xi)(Gg)_i=g_i$, with its holomorphic
interpretation at zero. Its kernels are diagonal, independent
of $\eta$, and have absolute value at most one along each
straight integration segment.

For every such diagonal kernel $D_0$, at arbitrary radii,
\[
 A_1(\xi)D_0A_1(\xi')=0,\qquad
 A_1(\xi)D_0\partial_\eta A_1(\xi')=0.
\]
Thus two adjacent derivative factors in an expansion of $K^k$
vanish, including the term in which the left derivative strikes
the right matrix. Each remaining word contains at most
$p_k=\lceil k/2\rceil$ parameter derivatives. Let a smaller
parameter neighborhood have distance $\Delta>0$ from the
boundary of the larger one. Split that loss evenly among the
at most $p_k$ derivatives and use Cauchy's integral estimate at
each differentiation. The ordered radial integration simplex
has volume $|\xi|^{k+1}/(k+1)!$; this follows by parametrizing
each nested segment by its ratio to the preceding complex
radius, so the same positive real simplex appears. Summing
the at most $2^k$ words and absorbing the bounded matrix and
source norms gives
\begin{equation}
 \|K^kGf_n\|_{|\xi|\leq R',\,\Omega'}
 \leq C_n^{k+1}\frac{(R')^{k+1}}{(k+1)!}
           \max\{1,p_k/\Delta\}^{p_k},\qquad R'<R .
 \label{eq:cq-actual-complex-axis-Picard-bound}
\end{equation}
The estimate applies to complete derivative operators, so every
Leibniz derivative of a coefficient is included.
Its $k$th root tends to zero: the factorial has root comparable
to $k$, whereas the derivative loss has root at most a constant
times $k^{1/2+1/(2k)}$. Hence
$\sum_{k\geq0}K^kGf_n$ converges normally to a holomorphic
solution of \eqref{eq:cq-actual-axis-linear-system}. Applying
the same estimate to the iterated homogeneous integral equation
proves uniqueness among bounded parameter-holomorphic solutions
on the real restricted radial interval. The source's actual
solution has that parameter regularity by Lemma 5.1; its
restriction therefore equals this holomorphic solution.

The source equations preserve even parity in the first four
components and odd parity in the last two, as stated and proved
by its uniqueness argument on page 49. This also follows by
reflecting the just-constructed integral equation and using the
same uniqueness. The first four convergent complex Taylor
series therefore contain only powers $\xi^{2j}=X^j$.
Their radius in $X$ is the square of a positive radius in
$\xi$, proving analytic $X$ germs for $\phi_n,U_n,K_n,\Pi_n$.
Formula (5.2) and the analytic radial average prove the same
for $V_n/X$. This closes the induction for the actual profiles,
including their full parameter derivatives.

We next apply this germ statement to the summed source field.
On a neighborhood of each fixed preterminal $(z,t)$, $q(z,t)$
has a positive lower bound. Source equation (10.1), the
cutoff-summed slow field and equation (9.21) consequently have
only finitely many active expansion coefficients and correction
stages there. The summation cutoffs are functions of $q(z,t)$,
independent of $r$; their smooth derivatives therefore multiply
the analytic radial germs without changing radial analyticity.
All active annular representatives vanish on an initial radial
interval, by the source support clauses, and the finite
intersection of those intervals still has positive length.
The locally summed field near the axis is thus the actual
axisymmetric slow field, with coefficients analytic in $r^2$.
The same statement holds after any fixed $z,t$ derivatives,
since the active sum is finite and each germ is jointly analytic
in $(X,\eta)$ before the smooth $q$-dependent factors.

The final spatial cutoff $\chi_x$ is exactly one for
$r<r_0/2$, $|z|<z_0/2$ in the original viscosity-one coordinates.
The temporal cutoff depends only on $t$. Hence throughout this
central axial slab, for every $t<1$, the actual localized
velocity and pressure have analytic radial germs. If the
temporal cutoff vanishes on the local time interval the fields
are zero and the conclusion is immediate; on its support the
preceding local source construction applies. The force before
time one is the full momentum residual, including every cutoff
derivative. It also has analytic radial germs there: the
Cartesian velocity can be written
\[
 u_1=\alpha(\sigma,z,t)x_1-\beta(\sigma,z,t)x_2,\quad
 u_2=\alpha(\sigma,z,t)x_2+\beta(\sigma,z,t)x_1,\quad
 u_3=\gamma(\sigma,z,t)
\]
with radial-analytic coefficients and all needed parameter
derivatives. Cartesian differentiation, multiplication and
addition preserve this form of radial analyticity. No
nonanalytic spatial cutoff is present on the stated slab.

Physical scaling $x=\sqrt\nu\,y$ takes the radial argument
to $\sigma/\nu$ and the axial argument to $z/\sqrt\nu$.
It therefore preserves analytic radial germs for every
$\nu>0$. In particular the actual $B,P,Q,F$ are analytic at
$\sigma=0$ for
\begin{equation}
 |z|<\sqrt\nu\,z_0/2,\qquad 0\leq t<1.
 \label{eq:cq-actual-central-analytic-domain}
\end{equation}
Their angular projection formulas make this explicit:
$B=2\sigma\beta$, $P=4\sigma^2\alpha\beta$,
$Q=2\sigma\gamma\beta$ in the local axisymmetric region;
the force torque has the identical radial factor
$F=2\sigma\beta_f$. These formulas retain the complete
locally summed fields and their coefficients. Thus the
radial entrance theorem applies to the actual untruncated
four fields on \eqref{eq:cq-actual-central-analytic-domain}:
\[
 \mathcal F_B^{\rm rad},\quad\mathcal F_P^{\rm rad},\quad
 \mathcal F_Q^{\rm rad},\quad\mathcal F_F^{\rm rad}
           \in\mathcal A_{\rm Sch}\subset\Xi\cqE .
\]
The compact radial support was already proved, so the endpoint
analyticity just established is all that was needed for the
order-one sufficient criterion. At other axial points, the
source only prescribes a smooth spatial cutoff; the previously
constructed full-input packet ingress retains those points
without imposing an unsupported analytic cutoff.

\subsection{A fixed entire spectral value of the actual concentrating field}

At $z=0$, $q=\tau$, $\eta=0$ in the original source coordinates.
The source's equation (B.3) gives $\phi_0(0,0)=1$; Lemma 5.1
gives $\phi_n(0,0)=0$ for every $n\geq1$.
Its locally finite sum, the vanishing of all annular corrections
near the axis, and the final cutoffs equal to one near the
specified terminal central point yield the exact axial identity
\begin{equation}
 \frac{u_{\nu,\theta}}r\bigg|_{\sigma=0,z=0,t=1-\tau}
          =\frac{\tau^{-1-h}}C,\qquad
 \partial_\sigma B(0,0,1-\tau)=\frac2C\tau^{-1-h}.
 \label{eq:cq-actual-axis-angular-coefficient}
\end{equation}
The first coefficient has the original amplitude $C>1$ from
the source and the original power
$q^{-A-1/2}=q^{-1-h}$. The second identity follows from
$B=r u_\theta=2\sigma(u_\theta/r)$.
Under viscosity scaling $u_\nu=\sqrt\nu\,u_*(x/\sqrt\nu,t)$,
both the numerator $u_\theta$ and radius $r$ have the same
factor $\sqrt\nu$, which cancels in this ratio. Thus no
viscosity factor is absent from the displayed coefficient.

For the direct arithmetic Mellin profile this gives the genuine
uncanceled pole
\[
 \operatorname{Res}_{w=-1}W_h(w)\widehat B(w,0,1-\tau)
       =\frac{2^{h+1}-1}{8C}\tau^{-1-h}.
\]
Here $W_h(-1)=(2^{h+1}-1)/16$ follows from
$\zeta(-1)=-1/12$ and the two retained dyadic factors.
The pole is compatible with the entire radial ingress because
the two precise maps have different Mellin arguments and
different moment-canceling dilation coefficients.

For the latter, set the fixed original spectral point
$s_0=-5i/2$, so that $w=1/2-is_0=-2$.
Formula \eqref{eq:cq-angular-radial-Taylor-data} gives
\[
 \operatorname{Res}_{w=-2}M_B^{\rm rad}(w,0,1-\tau)
   =-\frac{(1-2^{-3/2})(2^{h+1}-1)}C\,\tau^{-1-h}.
\]
The completed zeta function determines the exact canceled
zero coefficient:
\[
 \xi(-2)=\xi(3)=\frac{3\zeta(3)}{2\pi},\qquad
 \Gamma(w/2)=-\frac2{w+2}+O(1),\qquad
 \zeta'(-2)=-\frac{\zeta(3)}{4\pi^2}.
\]
For the first equality evaluate the retained completion at
$3$, using $\Gamma(3/2)=\sqrt\pi/2$, and use the functional
equation $\xi(-2)=\xi(3)$. For the second, the gamma recurrence
gives the residue $-1$ at its argument $-1$, and the change
of argument $w/2$ supplies the factor two. Multiplication by
$w(w-1)\pi^{-w/2}/2$, whose value at $-2$ is $3\pi$, shows
$\xi(-2)=-6\pi\zeta'(-2)$, proving the third equality.

Multiplying the simple zeta zero by the exact Mellin pole
therefore gives the actual entire trace value
\begin{equation}
 \mathcal F_B^{\rm rad}(s_0,0,1-\tau)
 =\frac{\zeta(3)}{4\pi^2C}
       (1-2^{-3/2})(2^{h+1}-1)\tau^{-1-h}.
 \label{eq:cq-actual-entire-trace-fixed-value}
\end{equation}
The function here is the untruncated actual radial source trace
already proved to lie in $\Xi\cqE$ for each preterminal time.
Writing its unique entire quotient as
$\mathcal F_B^{\rm rad}=\Xi g_B^{\rm rad}$ and retaining
$\Xi(s_0)=\xi(3)=3\zeta(3)/(2\pi)$ gives
\begin{equation}
 g_B^{\rm rad}(s_0,0,1-\tau)
 =\frac{(1-2^{-3/2})(2^{h+1}-1)}{6\pi C}\tau^{-1-h}.
 \label{eq:cq-actual-entire-quotient-fixed-value}
\end{equation}
All factors are strictly positive, so both displayed values
diverge with exactly the source exponent. This is a
fixed-spectral-point consequence for actual order-one entire
functions, not an uninverted trace claim at a moving sample.

Its precise relation to the source quotient retains the original
relation space. Since $\Xi(s_0)\ne0$, the function
$e(s)=\Xi(s)/\Xi(s_0)$ belongs to
$\mathbb C[s]\Xi\subseteq\cqN$ and satisfies $e(s_0)=1$.
For any representative $f$ and any prescribed scalar $b$,
$f+(b-f(s_0))e$ has the same source quotient class and
value $b$ at $s_0$. Thus the representative relation
$\{([f]_{\cqQ},f(s_0)):f\in\cqE\}$ has full scalar fibres.
This proves the exact information lost by imposing the source
quotient on the now-constructed growing trace value.
The full trace itself and its inverse radial input remain
specified by the preceding injective maps and kernels;
the calculation does not decide the distinct membership of
$\Xi'$ in $\cqN$.
